\documentclass[10pt,a4paper]{amsart}
\usepackage[margin=19mm]{geometry}
\usepackage{amsmath,amssymb,mathtools}
\usepackage[scr=rsfs,cal=euler]{mathalfa}
\usepackage{xfrac}
\usepackage[unicode,hidelinks,bookmarks=false]{hyperref}
\newcommand{\ie}{i.e.\ }
\newcommand{\cf}{cf.\ }
\newcommand{\resp}{respectively\ }
\newcommand{\Subsection}{\subsection}
\providecommand{\nobreakdash}{\nobreak\hbox{-}}
\setlength{\emergencystretch}{2em}
\multlinegap0pt
\newtheorem{defn}{Definition}
\newtheorem{ass}{Assumption}
\newtheorem{thm}{Theorem}
\newtheorem{pf}{Pf}
\title{Turnpike properties in nonlinear system identification}
\author{Julian D. Schiller, Matthias A. M\"uller}
\date{Source: arXiv:2609.02071v1 (CC BY 4.0)}
\begin{document}
\maketitle

\noindent\emph{Source and licence.} Turnpike properties in nonlinear system identification. Version arXiv:2609.02071v1 (CC BY 4.0), distributed by arXiv under the Creative Commons Attribution 4.0 International licence, http://creativecommons.org/licenses/by/4.0/, as recorded in the arXiv licence metadata for this exact version and shown on the abstract page https://arxiv.org/abs/2609.02071v1. Full licence notice: \texttt{licenses/arxiv-2609.02071-CC-BY-4.0.txt}.\par\smallskip

\noindent\textit{Editorial scope.} Counted unit: the article's thm thm:coerc (source0.tex lines 638--646). The article's complete original proof of it is delivered with it (source0.tex lines 648--683, one \texttt{proof} environment of 36 lines). No mathematics is written, repaired or re-derived: the statement and the proof are the article's own lines, in the article's own notation, and no retained line is substituted or re-flowed.  Retained uncounted as the source-local context the proof needs: lines 137--144; lines 164--166; lines 228--230; lines 234--240; lines 575--585; lines 590--596.  Nothing was omitted from any retained span, so the package carries no omission record.  No retained line contains a citation, so the wrapper carries no bibliography and no bibliography entry.  No retained line contains a figure environment, an image-inclusion command or drawing code, so no image is delivered and the package declares no asset dependency.  Editorial formatting only, in addition: an AMS \texttt{amsart} preamble that restates the article's own declarations and macros used by the retained lines, namely the environments \texttt{defn}, \texttt{ass}, \texttt{thm}, \texttt{pf} on separate counters, together with the standard packages those lines need.  The retained environments print in this document as Definition 1, Ass 1, Theorem 1, Pf 1 in order of appearance, whereas the article prints the same items with the numbering of its own full document.  The complete native source of the article as distributed in the arXiv e-print for this exact version is bundled unchanged as \texttt{sources/209171/source0.tex}.
	\begin{align}
		&\min\limits_{x_0,\theta} J(x_0,\theta) \label{eq:NLP_min}\\
		\text{s.t. \quad }
		& x_{j+1} = f(x_j,u_j^\mathrm{d};\theta), \ j\in[0,N], \label{eq:NLP_f}\\
		& y_{j} = h(x_{j},u_j^\mathrm{d};\theta), \ j\in[0,N],\label{eq:NLP_h}\\
		& x_{j+1}\in\mathcal{X}, \ y_{j} \in \mathcal{Y},\ j\in[0,N] \label{eq:NLP_X}\\
		& x_0\in\mathcal{X}_0, \ \theta \in \Theta. \label{eq:NLP_Theta}
	\end{align}

\begin{equation}\label{eq:NLP_case_1}
	J(\bar{x},\theta') = V'(\bar{x}) := \min_{\theta}\{ J(\bar{x},\theta)\ |\  \text{\eqref{eq:NLP_f}--\eqref{eq:NLP_Theta}}\} < \infty.
\end{equation}

\begin{equation}\label{eq:B_def}
	M(N):=\frac{1}{N+1}\sum_{j=0}^N\beta_j.
\end{equation}

\begin{defn}[Coercivity]\label{def:coerc}
	The value function in~\eqref{eq:NLP_case_1} is \emph{coercive} if there exist $\gamma,\alpha_\mathrm{c}\in\mathcal{K}_\infty$ and $C_\mathrm{c}\geq0$ such that for all $N\in\mathbb{N}$ and $\bar{x}\in\mathcal{X}_0$, it holds that
	\begin{equation}\label{eq:coerc}
		V'(\bar{x}) \geq V^* + M(N)\left(\gamma\left( \sum_{j=0}^{N} \alpha_\mathrm{c}(\|y'_j - y^*_j\|)\right) - C_\mathrm{c}\right)
	\end{equation}
	for all $\theta'\in\mathcal{S}'(\bar{x})$ and all $(x_0^*,\theta^*)\in \mathcal{S}^*(\bar{x},\theta')$, where $y_j' = y_j(\bar{x},\theta',U^\mathrm{d})$ and $y_j^* = y_j(x_0^*,\theta^*,U^\mathrm{d})$, $j\in[0,N]$.
\end{defn}

\begin{ass}[Uniform strict convexity]\label{ass:convex_lr}
	The function $l$ is uniformly strictly convex with modulus $\phi\in\mathcal{K}_\infty$, i.e., it holds that
	\begin{equation*}
		\begin{split}
			l(ty_1+(1-t)y_2,y^\mathrm{d})
			\leq&\ tl(y_1,y^\mathrm{d})+(1-t)l(y_2,y^\mathrm{d})\nonumber\\
			&\ -t(1-t)\phi(\|y_1-y_2\|)
		\end{split}
	\end{equation*}
	for all $y_1,y_2\in\mathcal{Y}$ uniformly for all $y^\mathrm{d}\in\mathcal{Y}^\mathrm{d}$ and $t\in[0,1]$.
\end{ass}

\begin{ass}[Optimality condition]\label{ass:optimality}
	There exists a constant $K\geq0$ such that, for all $N\in\mathbb{N}$ and $\bar{x}\in\mathcal{X}_0$, it holds that
	\begin{equation}\label{eq:cond_V_Jt}
		V^* - J(t) \leq M(N)K, \ \ \ J(t) := \sum_{j=0}^N\beta_jl(y_j(t),y^\mathrm{d}_j), \ \ t\in[0,1]
	\end{equation}
	for all $\theta'\in\mathcal{S}'(\bar{x})$ and all $(x_0^*,\theta^*)\in\mathcal{S}^*(\bar{x},\theta')$, where $y_j(t) = ty_j(\bar{x},\theta',U^\mathrm{d}) + (1-t)y_j(x_0^*,\theta^*,U^\mathrm{d})$, $j\in[0,N]$, $t\in[0,1]$.
\end{ass}

\begin{thm}\label{thm:coerc}	
	Let Assumptions~\ref{ass:convex_lr} and~\ref{ass:optimality} be satisfied.
	Suppose that there exists a constant $\underline{\beta}>0$ such that%
	\begin{equation}\label{eq:beta_lower_bound}
		\beta_j \geq \underline{\beta}M(N), \quad j\in[0,N]
	\end{equation}
	uniformly for all $N\in\mathbb{N}$ with $M(N)$ from~\eqref{eq:B_def}.
	Then, the value function in~\eqref{eq:NLP_case_1} is coercive in the sense of Definition~\ref{def:coerc}.
\end{thm}

\begin{pf}
	Define ${l}_\Delta(y_1,y_2,y^\mathrm{d}):= l(y_1,y^\mathrm{d}) - l(y_2,y^\mathrm{d})$ for $y_1,y_2,y^\mathrm{d}\in\mathbb{R}^p$.
	Consider arbitrary $y_1,y_2\in\mathcal{Y}$ and $y^\mathrm{d}\in\mathcal{Y}^\mathrm{d}$ and introduce $y(t) := ty_1 + (1-t)y_2$, $t\in[0,1]$.
	From Assumption~\ref{ass:convex_lr}, we have
	\begin{align*}
		&l_\Delta(y(t),y_2,y^\mathrm{d}) = l(y(t),y^\mathrm{d}) - l(y_2,y^\mathrm{d}) \\
		&\leq tl(y_1,y^\mathrm{d})+(1-t)l(y_2,y^\mathrm{d}) \\
		&\quad - t(1-t)\phi(\|y_1-y_2\|) - l(y_2,y^\mathrm{d})\\
		&= tl(y_1,y^\mathrm{d})-tl(y_2,y^\mathrm{d}) - t(1-t)\phi(\|y_1-y_2\|)\\
		&= tl_\Delta(y_1,y_2,y^\mathrm{d}) - t(1-t)\phi(\|y_1-y_2\|).
	\end{align*}
	Rearranging and dividing by $t\in(0,1)$ yields
	\begin{equation*}
		l_\Delta(y_1,y_2,y^\mathrm{d})  \geq \frac{1}{t}l_\Delta(y(t),y_2,y^\mathrm{d}) + (1-t)\phi(\|y_1-y_2\|).
	\end{equation*}
	Using the definition of $l_\Delta$ leads to
	\begin{equation}
		\begin{split}\label{eq:convexity_res1}
		l(y_1,y^\mathrm{d}) - l(y_2,y^\mathrm{d})
		\geq &\ \frac{1}{t}l(y(t),y^\mathrm{d}) - \frac{1}{t}l(y_2,y^\mathrm{d}) \\
		&+(1-t)\phi(\|y_1-y_2\|).
		\end{split}
	\end{equation}

	Consider any $N\in\mathbb{N}$, $\bar{x}\in\mathcal{X}_0$, $\theta'\in\mathcal{S}'(\bar{x})$, and $(x_0^*,{\theta}^*)\in\mathcal{S}^*(\bar{x},\theta')$, generating the outputs $y'_j:=y_j(\bar{x},{\theta}',U^\mathrm{d})$ and ${y}_j^*:=y_j({x}^*_0,{\theta}^*,U^\mathrm{d})$, $j\in[0,N]$. Moreover, for each $j\in[0,N]$, let ${y}_j(t) := ty'_j + (1-t)y^*_j$, $t\in[0,1]$. From~\eqref{eq:convexity_res1} and Assumption~\ref{ass:optimality}, we can immediately conclude that
	\begin{equation}
		V'(\bar{x}) - V^* \geq -M(N)\frac{K}{t} + (1-t)\sum_{j=0}^N\beta_j\phi(\|y'_j-{y}^*_j\|)\label{eq:proof_convexity}
	\end{equation}
	with $t\in(0,1)$.
	Now arbitrarily fix $t\in(0,1)$ and let $c:=1-t\in(0,1)$ and $\bar{K}:=K/t\geq0$. Consequently, from~\eqref{eq:proof_convexity} and the lower bound in~\eqref{eq:beta_lower_bound}, we can infer that
	\begin{equation*}
		V'(\bar{x}) - V^* \geq -M(N)\bar{K} +  M(N)\underline{\beta}c\sum_{j=0}^N\phi(\|y'_j-{y}^*_j\|)\label{eq:proof_convexity_2}
	\end{equation*}
	uniformly for $N\in\mathbb{N}$, $\bar{x}\in\mathcal{X}_0$, $\theta'\in\mathcal{S}'(\bar{x})$, and $(x_0^*,{\theta}^*)\in\mathcal{S}^*(\bar{x},\theta')$.
	Setting $\gamma(s) = \underline{\beta}c\cdot s$, $\alpha_\mathrm{c} = \phi$, and $C_\mathrm{c}= \bar{K}$ establishes the coercivity property from Definition~\ref{def:coerc}, thus completing the proof.
\end{pf}

\end{document}
